\chapter{Background}
\label{chap:background}

\section{Transformers and the Residual Stream}

A transformer language model processes a sequence by repeatedly updating a vector representation for every token. Each block contains an attention sublayer, which moves information between token positions, and a multilayer perceptron (MLP), which performs a position-wise nonlinear transformation \cite{vaswani2017attention}. Both communicate through the \emph{residual stream}: instead of replacing the current representation, each sublayer adds its output to it. If $x_{\ell,t}\in\mathbf{R}^{d}$ is the residual state at layer $\ell$ and token position $t$, a simplified update is

\begin{align}
x'_{\ell,t} &= x_{\ell,t}+\operatorname{Attention}_{\ell}(x_{\ell,1:T})_t, \\
x_{\ell+1,t} &= x'_{\ell,t}+\operatorname{MLP}_{\ell}(x'_{\ell,t}).
\end{align}

The residual stream is therefore a shared communication channel: attention heads and MLPs read from it, add information, and pass the combined state to later layers. Because output logits are ultimately computed from the final residual state, intermediate residual activations are a natural location for studying information that contributes to a response.

This project examines the Solver's residual-stream state after it has read the Critic feedback and immediately before it generates its decision. We cache the activation at the final input token, producing one fixed-size vector per interaction. This captures the state from which generation begins, but not every token-level computation; later chapters treat that limitation explicitly.

\section{Mechanistic Interpretability}

Mechanistic interpretability seeks explanations of model behavior in terms of internal representations and computations rather than input--output correlations alone. Here, the target behavior is whether the Solver adopts the Critic-supported answer or retains its original answer. Three questions must remain distinct:

\begin{enumerate}
    \item \textbf{Decodability:} can the outcome be predicted from the activation?
    \item \textbf{Interpretability:} can the feature be given a stable, human-readable description?
    \item \textbf{Causal relevance:} does changing the representation alter behavior in the direction predicted by that description?
\end{enumerate}

A representation may satisfy one criterion without satisfying the others. A probe might decode answer position even if the Solver does not use that signal to evaluate feedback. Conversely, an intervention might change generation by damaging the model rather than isolating a meaningful concept. We therefore combine semantic inspection, held-out prediction, and controlled intervention.

\subsection{Superposition and Polysemanticity}

Activation dimensions are not generally expected to correspond one-to-one with concepts. Neural networks can represent more useful variables than they have dimensions by encoding them as directions in activation space, an arrangement commonly described as \emph{superposition} \cite{elhage2022superposition}. One consequence is \emph{polysemanticity}: the same neuron or coordinate may participate in several unrelated patterns, while one concept may be distributed across many coordinates.

Superposition makes coordinate-level inspection difficult. A large value in one residual dimension need not have a stable meaning across contexts, and zeroing it can affect several overlapping representations. Interpretability methods therefore search for directions, subspaces, or sparse decompositions that align more closely with meaningful variables. Sparse autoencoders are one such method, but their latents remain candidate features rather than guaranteed semantic units.

\subsection{Probing and Its Limits}

A probe is a supervised model trained to predict an external label from an internal activation. Linear probes restrict the decoder's complexity and expose a weighted direction in representation space \cite{alain2016probes}. In this study, probes test whether SAE latents distinguish Adoption from Retention, and their coefficients rank candidate features for deeper analysis.

Probe accuracy does not show that the model uses the decoded information. A probe can exploit incidental correlations, excess capacity, or dataset shortcuts, and its coefficients can be unstable when features are correlated. Control tasks and selectivity measurements help distinguish target structure from memorization \cite{hewitt2019control}. We use held-out question families, control probes, and repeated random-feature draws, treating feature ranking as discovery rather than a final interpretation.

\subsection{Causal Methods}

Causal methods intervene on internal activations and observe the output. \emph{Ablation} removes or suppresses a component to test whether it is necessary. \emph{Activation patching} replaces an activation with one from a comparison run to localize a behaviorally relevant difference \cite{heimersheim2024patching}. \emph{Steering} adds or subtracts a direction to test whether moving the representation toward a hypothesized state changes generation.

These methods provide stronger evidence than observation alone, but large edits may move activations out of distribution, damage generation, or change multiple correlated variables. An output flip is therefore not automatically evidence for the intended mechanism. This project compares targeted edits with equal-size random edits, records edit norms, measures SAE reconstruction without steering, and separates valid Adoption--Retention switches from broken outputs.

\section{Sparse Autoencoders}

An SAE learns a sparse, usually overcomplete representation of dense activations. For $x\in\mathbf{R}^{d}$, the encoder produces $z\in\mathbf{R}^{m}$, often with $m>d$, and the decoder reconstructs the activation:

\begin{align}
z &= f\!\left(W_{\mathrm{enc}}(x-b_{\mathrm{dec}})+b_{\mathrm{enc}}\right), \\
\hat{x} &= W_{\mathrm{dec}}z+b_{\mathrm{dec}}.
\end{align}

Decoder columns define feature directions in activation space, while entries of $z$ indicate feature strength. Training balances reconstruction with sparsity. One common objective is

\begin{equation}
\mathcal{L}=\lVert x-\hat{x}\rVert_2^2+\lambda\lVert z\rVert_1,
\end{equation}

whereas a TopK SAE retains only the $k$ largest latent activations. Reconstruction encourages information preservation; sparsity encourages each example to use a small subset of the dictionary.

SAEs are attractive because sparsity makes individual latents inspectable and the decoder maps a latent edit back into the residual stream \cite{cunningham2023sparse}. A latent may appear interpretable when its highest-activating examples share a pattern. That is insufficient by itself: feature splitting can distribute one concept across latents, while feature absorption can make a seemingly coherent latent fail where it should activate \cite{chanin2024absorption}. Reconstruction and sparsity are thus diagnostics, not direct measures of semantic or causal faithfulness.

We compare a broadly trained pretrained SAE, custom SAEs trained on Solver--Critic activations, and adapted variants. The pretrained model comes from Gemma Scope, a collection of open SAEs trained across layers and sites in Gemma models \cite{lieberum2024gemmascope}. Gemma Scope offers a reusable general-purpose dictionary, whereas the custom SAEs are fitted to the activation distribution created by this project's feedback interactions. Adapted variants modify an existing decoder while preserving part of its learned structure. These comparisons test whether domain alignment improves reconstruction, prediction, feature coherence, or causal control rather than assuming that one source is inherently more interpretable.

\section{LLM Feedback and Critique}

Solver--Critic systems are a form of multi-agent inference: one model proposes an answer, another supplies an evaluation or alternative, and the Solver responds again with access to that feedback. Debate and critique can improve performance when agents contribute complementary evidence \cite{du2024debate}. The central behavioral question is nevertheless not simply whether the final answer changed, but whether the Solver adopted the Critic-supported answer, retained its own, or already agreed.

Revision may arise from mechanisms other than evaluating the feedback. Models may defer to assertion, follow presentation-order cues, imitate confidence, or align with an external belief. Sycophancy research shows that agreement can conflict with truthfulness \cite{sharma2024sycophancy}; self-correction research shows that reconsideration without informative external feedback does not reliably improve reasoning \cite{huang2024selfcorrect,kamoi2024selfcorrection}. Iterative methods such as Self-Refine demonstrate that structured feedback can improve outputs in some settings \cite{madaan2023selfrefine}, while knowledge-conflict studies show that models vary between adapting to contextual evidence and preserving parametric knowledge \cite{xie2024chameleon}. Uptake must therefore be evaluated separately from feedback correctness.

Our taxonomy isolates this transition. Adoption means moving from the original answer to the Critic-supported candidate; Retention means preserving the original answer; Agreement covers substantively equivalent candidates; and unresolved outputs are excluded from the binary comparison. This target gives feature analysis a defined phenomenon to explain, while stratification by correctness asks whether uptake was helpful or harmful.
